\subsection{Operating Theatre Scheduling Model\label{sec:ot}}
\info[author=Joao]{Depois preciso ``sincronizar'' essa seção com o todo que estamos produzindo, i.e., essa etapa de alocação vem depois das etapas de gerenciamento de filas e etc.}
We modeled a deterministic MILP\info[author=Joao]{Depois preciso garantir a introdução dessa sigla em seções anteriores.} formulation for the Operating Theatre Scheduling Problem (OTSP), which focuses on the tactical allocation of operating theatre blocks to medical subspecialities.

To formulate the OTSP, we use the Table~\ref{tab:symb} notation.
Note that the variable $\blockvar$ is a variable that indicates whether the theatre $\roomidx$ in the block $\blockidx$ is allocated to the subspeciality $\specialtyidx$.\unsure[author=Joao]{Isso bate com o comentario que ja fiz la em cima: preciso sincronizar no texto o que é bloco e o que é alocação!}
Therefore, throughout the text, we will use the term \textit{allocation} to refer to an allocated speciality.

\begin{table}
    \centering
    \caption{\label{tab:symb}Table of Symbols.}
    \scriptsize
    \begin{tabular}{cl}
    \toprule
    \textbf{Set}                    & \textbf{Description}                                                                       \\
    \midrule
        $\roomset$                       & Operating theatres (both elective and outpatient)                       \\
        $\planninghorizonset$            & Planning horizon days, such that $\planninghorizonset = \{1, 2, ..., \days\}$  \\
        $\timeperiodset$                 & Time periods (morning or afternoon), such that $\timeperiodset = \{\text{morning}, \text{afternoon}\}$          \\
        $\blockset$                      & Blocks, such that $\blockset = \planninghorizonset \times \timeperiodset$     \\
        $\specialtyset$                  & Surgical subspecialities                                                                \\
        $\specialtymset$                 & Surgical subspecialities that must have two sequential blocks allocated,\\
                                         & such that $\specialtymset \subseteq \specialtyset$          \\
        $\specialtycset$                 & Circumstantial surgical subspecialities (do not occur every month),                      \\
                                         & but have high priority, such that $\specialtycset \subseteq \specialtyset$              \\
        $\specialtyprset$                & Priority surgical subspecialities, such that $\specialtyprset \subseteq \specialtyset$  \\
        $\specialtytxset$                & Organ transplant surgical subspecialities, such that $\specialtytxset \subseteq \specialtyset$                                                                                                              \\
    \midrule
    \textbf{Index}                           & \textbf{Description}                                                                  \\
    \midrule
        $\roomidx$                            & Operating theatre, with $\roomidx \in \roomset$                                       \\
        $\planninghorizonidx$                 & Planning horizon day, with $\planninghorizonidx \in \planninghorizonset$ \\
        $\timeperiodidx$                      & Time period (morning or afternoon), with $\timeperiodidx \in \timeperiodset$           \\
        $\blockidx$                           & Surgery block, with $\blockidx \in \blockset$                                    \\
        $\specialtyidx$                       & Surgical subspeciality, with $\specialtyidx \in \specialtyset$                      \\
    \midrule
    \textbf{Parameter}                        & \textbf{Description}                                                                  \\
    \midrule
        $\allowedrooms_{\roomidx\blockidx\specialtyidx}$   & 1 if theatre $\roomidx \in \roomset$ in block $\blockidx \in \blockset$ can              \\
                                                           & be used by subspeciality $\specialtyidx \in \specialtyset$; 0 otherwise          \\
        $\demand_{\specialtyidx}$                          & Number of surgeries demanded for subspeciality $\specialtyidx \in \specialtyset$                       \\
        $\specteams_{\blockidx\specialtyidx}$              & Number of teams from subspeciality $\specialtyidx \in \specialtyset$ available in block  $\blockidx \in \blockset$ \\
        $\numbersanests_{\blockidx}$                       & Number of anesthesiologists available in block $\blockidx \in \blockset$                         \\
        $\needanest_{\specialtyidx}$                       & 1 if subspeciality $\specialtyidx \in \specialtyset$ requires an anesthesiologist; 0 otherwise         \\    
        $\mindeficit_{\specialtyidx}$                      & Minimum number of blocks required for subspeciality $\specialtyidx \in \specialtyset$            \\
        $\ub$                                              & Upper bound on the number of blocks to allocate, relative to each subspeciality $\specialtyidx \in \specialtyset$   \\
    \bottomrule
    \textbf{Variable}                         & \textbf{Description}                                                                  \\
    \midrule
        $\blockvar$                                        & 1 if theatre $\roomidx \in \roomset$ in block $\blockidx \in \blockset$ is allocated to subspeciality $\specialtyidx \in \specialtyset$; 0 otherwise                \\
        $\precvar$                                         & 1 if theatre $\roomidx \in \roomset$ on day $\planninghorizonidx \in \planninghorizonset$ is allocated exclusively to subspeciality $\specialtyidx \in \specialtyset$;\\
                                                           & 0 otherwise                                  \\
        $\integervar$                                      & Number of transplants to be performed in block $\blockidx \in \blockset$ by subspeciality $\specialtyidx \in \specialtyset$ \\
        $\balancevar$                                      & Minimum proportion of allocations per demand among all subspecialities   \\
    \bottomrule
    \end{tabular}
\end{table}

The Objective Functions (OFs) are organised according to the priority levels established by the hospital's guidelines.
To reflect these priorities, we adopt a hierarchical optimisation strategy, as proposed in~\citet{Aringhieri2022Hierarchical}.
Each OF is optimised sequentially, while respecting the optimal values obtained in the previous stages.

Formally, the first OF is defined as
\begin{align}
\text{Max } \displaystyle F_1 & = \sum_{\roomidx \in \roomset} \sum_{\blockidx \in \blockset} \sum_{\specialtyidx \in \specialtyprset} \blockvar, \label{eq:f0}
\end{align}
which aims to maximise the number of allocations for priority subspecialities ($\specialtyprset$).
The second function is given by
\begin{align}
\text{Max } \displaystyle F_2 & = \sum_{\roomidx \in \roomset} \sum_{\blockidx \in \blockset} \sum_{\specialtyidx \in \specialtyset} \blockvar, \label{eq:f1}
\end{align}
which seeks to maximise the total number of allocations, i.e., to increase theatre occupancy.
The third objective function is defined as
\begin{align}
\text{Max } \displaystyle F_3 & = \balancevar, \label{eq:f2}
\end{align}
and its purpose is to promote balance in the distribution of allocations across subspecialities.
This function works in conjunction with a model constraint that enforces that, for each $\specialtyidx \in \specialtyset$, the ratio between the number of allocations and the corresponding demand must be greater than or equal to $\balancevar$.
Hence, OF~\eqref{eq:f2} maximises the lowest allocation-to-demand ratio among all subspecialities, ensuring a balanced allocation and avoiding over-concentration in a few subspecialities to the detriment of others.

Finally, the fourth OF is given by
\begin{align}
\text{Max } \displaystyle F_4 & = \sum_{\roomidx \in \roomset} \sum_{\planninghorizonidx \in \planninghorizonset} \sum_{\specialtyidx \in \specialtyset} \precvar, \label{eq:f3}
\end{align}
which aims to maximise the number of subspecialities allocated sequentially on the same day in the same theatre.
This is particularly important for improving the logistics of the surgical process, reducing the need for equipment and team transfers between theatres.

With this, we formulate the theatre planning problem at HC - Unicamp by the $\modelname$ model, incorporating these OFs and constraints based on the hospital's current operational framework.\change[author=Joao]{O MODELO VAI MUDAR! De acordo com o que vai vir como input do RQ}

{\footnotesize
\linespread{1}\selectfont
\begin{alignat}{4}
    (\modelname) \quad & \omit\rlap{Maximize OFs~\eqref{eq:f0}-\eqref{eq:f3}}\nonumber\\
    & \mbox{subject to} && \quad & \sum_{\roomidx \in \roomset} \sum_{\blockidx \in \blockset} \blockvar  & \geq \mindeficit_{\specialtyidx}                    &\qquad & \forall\, \specialtyidx \in \specialtyset  \label{eq:m2-2}\\
    &                  && \quad & \sum_{\roomidx \in \roomset} \sum_{\blockidx \in \blockset} \blockvar  & \leq (1 + \ub) \cdot \demand_{\specialtyidx}   &\qquad & \forall\, \specialtyidx \in \specialtyset \label{eq:m2-3}\\
    &                  && \quad & \sum_{\specialtyidx \in \specialtyset} \blockvar                       & \leq 1     &\qquad & \forall\, \roomidx \in \roomset,\, \blockidx \in \blockset \label{eq:m2-4}\\
    &                  && \quad & \sum_{\roomidx \in \roomset} \sum_{\specialtyidx \in \specialtyset} \needanest_{\specialtyidx} \cdot \blockvar   & \leq \numbersanests_{\blockidx}                    &\qquad & \forall\, \blockidx \in \blockset \label{eq:m2-5}\\
    &                  && \quad & \sum_{\roomidx \in \roomset} \blockvar   & \leq \specteams_{\blockidx\specialtyidx}   &\qquad & \forall\, \,\blockidx \in \blockset,\,\specialtyidx \in \specialtyset \label{eq:m2-6}\\
    &                  && \quad & \blockvar   & = 0    &\qquad & \forall\, \roomidx \in \roomset,\,\blockidx \in \blockset,\,\specialtyidx \in \specialtyset : \nonumber\\
    &                  &&       &                                                          &                 &\qquad &\allowedrooms_{\roomidx\blockidx\specialtyidx} = 0 \label{eq:m2-7}\\
    &                  && \quad & \blockvar + x_{\roomidx(\blockidx+1)\specialtyidx} - 1   & \leq \precvar   &\qquad & \forall\, \roomidx \in \roomset,\, \blockidx \in \blockset,\,\specialtyidx \in \specialtyset, \nonumber\\
    &                  &&       &                                                          &                 &\qquad & \planninghorizonidx \in \planninghorizonset :\,\text{ $\blockidx$ and $\blockidx+1$ are blocks} \nonumber\\
    &                  &&       &                                                          &                 &\qquad & \text{ of the same day $\planninghorizonidx$}\label{eq:m2-8}\\
    &                  && \quad & \blockvar + x_{\roomidx(\blockidx+1)\specialtyidx}       & \geq 2 \cdot \precvar   &\qquad & \forall\, \roomidx \in \roomset,\, \blockidx \in \blockset,\,\specialtyidx \in \specialtyset, \planninghorizonidx \in \planninghorizonset : \nonumber\\
    &                  &&       &                                                          &                 &\qquad & \planninghorizonidx \in \planninghorizonset : \text{ $\blockidx$ and $\blockidx+1$ are blocks} \nonumber\\
    &                  &&       &                                                          &                 &\qquad & \text{ of the same day $\planninghorizonidx$}\label{eq:m2-9}\\
    &                  && \quad & \blockvar       & = x_{\roomidx(\blockidx+1)\specialtyidx}   &\qquad & \forall\, \roomidx \in \roomset,\, \blockidx \in \blockmset,\,\specialtyidx \in \specialtymset,\, \nonumber\\
    &                  &&       &                                                          &                 &\qquad & \planninghorizonidx \in \planninghorizonset : \text{ $\blockidx$ and $\blockidx+1$ are blocks} \nonumber\\
    &                  &&       &                                                          &                 &\qquad & \text{ of the same day $\planninghorizonidx$}\label{eq:m2-10}\\
    &                  && \quad & \sum_{\roomidx \in \roomset} \blockvar   & = 2 \cdot \integervar  & \qquad & \forall\, \blockidx \in \blockset,\,\specialtyidx \in \specialtytxset  \label{eq:m2-11}\\
    &                  && \quad & \sum_{\roomidx \in \roomset} \sum_{\blockidx \in \blockset} \blockvar      & = \demand_{\specialtyidx}     &\qquad & \forall\,\specialtyidx \in \specialtycset  \label{eq:m2-12}\\    
    &                  && \quad & \sum_{\roomidx \in \roomset} \sum_{\blockidx \in \blockset} \frac{\blockvar}{\demand_{\specialtyidx}}   & \geq \balancevar     &\qquad & \forall\,\specialtyidx \in \specialtyset  \label{eq:m2-13}\\ 
    &                  &&       & \blockvar                           & \in \{0,1\}                 &        & \forall\, \roomidx \in \roomset,\, \blockidx \in \blockset,\, \specialtyidx \in \specialtyset\nonumber \\
    &                  &&       & \precvar                            & \in \{0,1\}                 &        & \forall\, \roomidx \in \roomset,\, \planninghorizonidx \in \planninghorizonset,\, \specialtyidx \in \specialtyset\nonumber \\
    &                  &&       & \integervar                         & \in \mathbb{Z_{+}}          &        & \forall\, \blockidx \in \blockset,\, \specialtyidx \in \specialtyset\nonumber\\
    &                  &&       & \balancevar                         & \in \mathbb{R_{+}}.         &        &\nonumber
\end{alignat}
}

Constraints~\eqref{eq:m2-2} and~\eqref{eq:m2-3} impose minimum and maximum bounds on the allocation of each subspeciality.
The parameter $\mindeficit_{\specialtyidx}$, typically set to low values (between 1 and 3), ensures that no subspeciality is left completely unattended.
The upper bound, governed by $\ub \in [0,1]$, restricts allocations to at most $(1+\ub)$ times the demand $\demand_{\specialtyidx}$, avoiding over-allocation beyond a reasonable limit.\change[author=Joao]{Essas restrições vão mudar, de acordo com os inputs do gerenciamento de filas}

Constraint~\eqref{eq:m2-4} ensures that each theatre and time block receives at most one subspeciality.
Constraint~\eqref{eq:m2-5} limits the number of allocations that require anaesthetists according to their availability $\numbersanests_{\blockidx}$ in each block.
Constraint~\eqref{eq:m2-6} restricts the allocation of each subspeciality based on the availability of $\specteams_{\blockidx\specialtyidx}$ medical teams representing it in block $\blockidx$.

Compatibility between theatre, block, and subspeciality is ensured by constraint~\eqref{eq:m2-7}, while constraints~\eqref{eq:m2-8} and~\eqref{eq:m2-9} together guarantee sequential allocation of blocks on the same day when $\precvar = 1$.
Subspecialities requiring two consecutive blocks are specifically addressed by constraint~\eqref{eq:m2-10}.

Constraint~\eqref{eq:m2-11} imposes that the number of allocations for transplants must be even, reflecting the need for two theatres for the procedure -- one for the donor and one for the recipient.
Finally, constraint~\eqref{eq:m2-12} requires full service provision for circumstantial subspecialities $\specialtycset$, and equation~\eqref{eq:m2-13} promotes proportional allocation balancing in conjunction with OF~\eqref{eq:f2}.